\begin{textsample}{POS Dim 7 – global\_south\_2024\_05 – Score 3.00 – global\_south\_2024\_05/108452445.txt}  \label{ex:f7_pos_065}
Being recognised as a"state"by other countries can support a region ' s aspiration to be seen as a legitimate political entity by the international community.
Espen Barth Eide, Norway ' s Minister of Foreign Affairs, said in a post on X :"Norway will recognize Palestine as a State.
A two-state solution is the only viable pathway to peace for both Israel and Palestine.
At this critical juncture, our recognition comes in support of the work towards a comprehensive plan for regional peace."Advertisement In a belligerent response, Israel has recalled its ambassadors from the three countries.
What does such a recognition mean, where does it sit with how the world views Palestine, and why does it matter?
We explain.
Firstly, what does it mean to be recognised as a state?
The Montevideo Convention on the Rights and Duties of States ( 1933 ), identified four conditions of a state :"a permanent population, defined territory, government, and capacity to enter into relations with other states". @ @ @ @ @ @ @ @ @ @ Law,"has long been the central organising idea in the international system".
While several regions and peoples have over the years sought to declare themselves as independent states, their formal recognition depends on how the rest of the world views them.
The United Nations has a broad criterion for accepting states as Members.
Article 4 of the UN Charter states :"Membership in the United Nations is open to all other peace-loving states which accept the obligations contained in the present Charter and, in the judgment of the Organization, are able and willing to carry out these obligations."Procedurally, admission to the UN as a Member State is granted by a two-thirds majority vote in the UN General Assembly.
However, the UNGA takes up the candidature only upon the recommendation of the UN Security Council.
The UNSC comprises five permanent members -- the United States, the United Kingdom, Russia, China, and France -- and 10 temporary member countries chosen on a rotational basis.
For the UNSC @ @ @ @ @ @ @ @ @ @ at least nine members in favour and no permanent members using their \textbf{veto}.
Essentially, it is the P5 who determine the fate of an issue in the UNSC.
What is the status of Palestine at the UN?
Currently, Palestine is a"Permanent Observer State"-- and not a"Member State"-- at the UN.
There is one other Permanent Observer State in the UN -- the Holy See, representing Vatican City.
Advertisement As a Permanent Observer State, Palestine is allowed to"participate in all of the Organization ' s proceedings, except for voting on draft \textbf{resolutions} and decisions in its main organs and bodies, from the Security Council to the General Assembly and its six main committees".
Palestine graduated to the status of"non-member Permanent Observer State"from having Observer status in 2012.
Mahmoud Abbas, President of the Palestinian Authority, had then hoped for the exercise to"breathe new life"into the peace process in the region.
Palestine has @ @ @ @ @ @ @ @ @ @ in the past, most recently in April this year.
In the UNSC, the United States, Israel ' s staunchest ally, had \textbf{vetoed} its admission.
Which countries recognise Palestine as a state?
Before the announcement by Norway, Ireland, and Spain, 143 of the UN ' s 193 members already recognised Palestine as a state.
Most of these countries are in Asia, Africa, and South America.
India accorded recognition in 1988.
Advertisement Recognition as a state lies at the heart of the Palestinian people ' s right to self-determination, and to decide their political future and government.
The 1947 United Nations Partition Plan for Palestine ( UNGA \textbf{Resolution} 181(II) ) proposed the establishment of a Jewish state, an Arab state, and for the city of Jerusalem to be administered by the UN as a corpus separatum ( separate body ).
This is also known as the ' two-state solution '.
Section F of the \textbf{resolution} said when either of the states have become independent as envisaged in the @ @ @ @ @ @ @ @ @ @ their \textbf{application} for admission to UN membership.
However, Palestinian leaders rejected the Plan, which they believed went against Arab interests.
The Arab-Israeli war broke out soon afterward, and Israel emerged as the winner.
In 1949, the proposal for its UN membership was tabled, and all P5 members except the UK ( which abstained ) agreed.
Advertisement What is the significance of the Norway-Ireland-Spain move?
When a state recognises another, it usually leads to the setting up of an embassy and posting of diplomatic officials in that country.
The Norwegian Foreign Minister has said that its representative office to the Palestinian Authority, which was opened in the West Bank in 1999, would become an embassy.
The move marks a cleavage in the West ' s backing of Israel in the Gaza war, which has divided peoples and countries around the world.
The US and UK have earlier spoken of an independent Palestine, but they have said it should only be the result of a negotiated settlement with Israel. @ @ @ @ @ @ @ @ @ @ said he has taken the decision"out of moral conviction, for a just cause and because it is the only way that the two states, Israel and Palestine, can live together in peace".
He has said that while"fighting the terrorist group Hamas is legitimate and necessary",?
Israel ' s Prime Minister Benjamin Netanyahu"is creating so much pain and so much destruction and so much rancour in Gaza and the rest of Palestine that the two-state solution is in danger".

% matched lemmas: application, resolution, veto
\end{textsample}
